%% Description of the Herning test rig dataset.
%% Covers dimensions, combinations of faults/bearings/conditions, EDA findings,
%% supported tests and splits; the test rig is not described here.
%% Build (from eda/herning/):   pdflatex herning && pdflatex herning

\documentclass[11pt,a4paper]{article}
\usepackage[T1]{fontenc}
\usepackage[utf8]{inputenc}
\usepackage[margin=2.4cm]{geometry}
\usepackage{amsmath}
\usepackage{booktabs}
\usepackage{graphicx}
\usepackage[hidelinks]{hyperref}

\title{Herning test rig}
\author{}
\date{\today}

\begin{document}
\maketitle

\section{Source}

\section{Dimensions of the time series}

\section{Combinations of faults, bearings and operating conditions}

\begin{table}[h]
\centering
\small
\begin{tabular}{lll}
\toprule
Factor & Levels & $n$ \\
\midrule
Bearing            & one & 1 \\
Lubricant          & one & 1 \\
Temperature target & 40--100\,\textdegree C in 5\,\textdegree C steps, ca.\ 305 sweeps each & 13 \\
Speed              & swept within each temperature target & \\
\bottomrule
\end{tabular}
\end{table}

\begin{itemize}
  \item One run; plus 5 cool-down sweeps at target 25\,\textdegree C. Sweeps at rpm target
        100 are the bearing standing still.
  \item One sweep: 0.2\,s at 2.5\,MHz (500,000 samples), taken ca.\ every 12\,s; channels AE
        and UL as 16-bit codes (D22, D25).
  \item With one lubricant, $\kappa$ is a fixed function of rpm and temperature, so it cannot be
        separated from the operating point (QUESTIONS.md Q1).
  \item Status \emph{validation}: rig unfinished, final data pending.
\end{itemize}
% TODO: rpm levels per temperature target, from herning_eda.ipynb.

\section{Findings from EDA}

\section{Tests the data support}

\section{Splits}

\end{document}
